%! Composition of Functions — Unit 1, Lesson 1.5 — Group Activity
\documentclass[10pt]{article}
\usepackage{precalculus-article}
\usepackage{precalculus-boxes}
\newcommand{\TallMath}[1]{$\displaystyle #1\rule[-1.4em]{0pt}{3.2em}$}

\begin{document}

\pageheader{Unit 1, Lesson 1.5}{Group Activity}
\namepartnerperiod

\vspace{0.08in}

\begin{scenariobox}[Exploring Composition of Functions]{sky}
\small
In this activity you will evaluate, construct, and decompose compositions of functions.
Work with your group on the tier your teacher assigns. Be ready to explain your reasoning.
\end{scenariobox}

\vspace{0.08in}

\begin{tcolorbox}[colback=white, colframe=black!40,
    title=\textbf{Tier R --- Remediate},
    fonttitle=\bfseries, arc=2mm, left=3mm, right=3mm, top=2mm, bottom=2mm]

\textbf{Given:} $f(x) = 2x + 5$ and $g(x) = x - 3$.

\begin{enumerate}[itemsep=10pt]
  \item Evaluate $f(g(0))$ step by step.

        \vspace{0.04in}
        $g(0) = $ \blank{3cm} \quad then $f\!\left(\rule{1.5cm}{0.4pt}\right) = $ \blank{4cm}

  \item Evaluate $f(g(2))$ step by step.

        \vspace{0.04in}
        $g(2) = $ \blank{3cm} \quad then $f\!\left(\rule{1.5cm}{0.4pt}\right) = $ \blank{4cm}

  \item Evaluate $g(f(0))$ step by step.

        \vspace{0.04in}
        $f(0) = $ \blank{3cm} \quad then $g\!\left(\rule{1.5cm}{0.4pt}\right) = $ \blank{4cm}

  \item Based on your answers, is $f(g(0)) = g(f(0))$? What does that tell you about composition?

        \vspace{0.04in}
        \writeline
\end{enumerate}
\end{tcolorbox}

\vspace{0.08in}

\begin{tcolorbox}[colback=white, colframe=black!40,
    title=\textbf{Tier A --- Apply},
    fonttitle=\bfseries, arc=2mm, left=3mm, right=3mm, top=2mm, bottom=2mm]

\textbf{Table of values:}

\vspace{0.04in}
\renewcommand{\arraystretch}{1.6}
\begin{tabular}{c|c|c}
$x$ & $f(x)$ & $g(x)$ \\ \hline
1 & 5 & 3 \\
2 & 8 & 1 \\
3 & 11 & 4 \\
4 & 14 & 2 \\
\end{tabular}

\begin{enumerate}[itemsep=8pt]
  \item Fill in the table for $f(g(x))$ and $g(f(x))$.

        \vspace{0.04in}
        \renewcommand{\arraystretch}{1.6}
        \begin{tabular}{c|c|c}
        $x$ & $f(g(x))$ & $g(f(x))$ \\ \hline
        1 & \blank{2.5cm} & \blank{2.5cm} \\
        2 & \blank{2.5cm} & \blank{2.5cm} \\
        3 & \blank{2.5cm} & \blank{2.5cm} \\
        4 & \blank{2.5cm} & \blank{2.5cm} \\
        \end{tabular}

  \item Now let $f(x) = x^2$ and $g(x) = 2x - 1$. Write an analytic formula for $f(g(x))$.
        Simplify completely.

        \vspace{0.04in}
        $f(g(x)) = $ \blank{9cm}

  \item What is the domain of $f \circ g$ if $f(x) = x^2$ and $g(x) = 2x-1$?

        \vspace{0.04in}
        Domain of $f \circ g$: \blank{8cm}
\end{enumerate}
\end{tcolorbox}

\vspace{0.08in}

\begin{tcolorbox}[colback=white, colframe=black!40,
    title=\textbf{Tier E --- Extend},
    fonttitle=\bfseries, arc=2mm, left=3mm, right=3mm, top=2mm, bottom=2mm]

\textbf{Given:} $h(x) = (3x + 2)^4$.

\begin{enumerate}[itemsep=10pt]
  \item Find one decomposition $h(x) = f(g(x))$. State $f$ and $g$, then verify.

        \vspace{0.04in}
        $g(x) = $ \blank{5cm} \qquad $f(u) = $ \blank{5cm}

        Verify: $f(g(x)) = $ \blank{7cm} $= h(x)$? \blank{2cm}

  \item Find a \textbf{different} valid decomposition $h(x) = p(q(x))$. State $p$ and $q$, then verify.

        \vspace{0.04in}
        $q(x) = $ \blank{5cm} \qquad $p(u) = $ \blank{5cm}

        Verify: $p(q(x)) = $ \blank{7cm} $= h(x)$? \blank{2cm}

  \item For your first decomposition, explain: is $h(x) = f(g(x))$ or $h(x) = g(f(x))$?
        Write out $g(f(x))$ and show whether it equals $h(x)$.

        \vspace{0.04in}
        $g(f(x)) = $ \blank{8cm}

        \vspace{0.04in}
        Does $g(f(x)) = h(x)$? \blank{2cm} \quad Explain: \writeline
\end{enumerate}
\end{tcolorbox}

\end{document}
